\documentclass[11pt]{article}
\usepackage[margin=1.1in]{geometry}
\usepackage{amsmath,amssymb,amsthm}
\usepackage{booktabs}
\usepackage{graphicx}
\usepackage[hidelinks]{hyperref}

\newtheorem{theorem}{Theorem}
\newtheorem{lemma}[theorem]{Lemma}
\newtheorem{proposition}[theorem]{Proposition}
\newtheorem{corollary}[theorem]{Corollary}
\newtheorem{conjecture}[theorem]{Conjecture}
\theoremstyle{definition}
\newtheorem{definition}[theorem]{Definition}
\theoremstyle{remark}
\newtheorem{remark}[theorem]{Remark}

\newcommand{\R}{\mathbb{R}}
\DeclareMathOperator{\asin}{arcsin}

\title{Arranging circles of radii $1,2,\dots,n$ around a central circle:\\
a Supnick TSP and certified finite optima}
\author{Maurizio Falconi\thanks{Independent researcher, Parma, Italy. \texttt{maurizio.falconi47@gmail.com}}}
<<<<<<< HEAD
\date{June 2026}
=======
\date{June 2026 --- v0.2 (submission)}
>>>>>>> 2ed20549af73671a647ea86607f0acf2700730eb

\begin{document}
\maketitle

\begin{abstract}
We study a discrete-geometric optimization problem: circles of radii $1,2,\dots,n$ are all externally tangent to a central circle, and the central radius $R$ is minimized over cyclic orders of the surrounding circles. We prove that the chain-ordering component is governed by a fixed Supnick/anti-Monge traveling-salesman order. For every $R$, the angular-separation matrix is symmetric anti-Monge, so Supnick's theorem gives one minimizing cyclic order, independent of $R$. This proves the conjectured ``pyramid'' order optimal whenever the corresponding chain necklace is geometrically realizable, and gives an unconditional lower bound in all cases.

Full geometric feasibility can fail because non-adjacent circle constraints are not captured by the chain equation; from $n=8$ the smallest circle can become a floating circle tangent only to the central circle. We formulate the full problem as a circular system of pairwise angular constraints, equivalently a simple temporal network, and certify global optima for $3\le n\le14$ using branch-and-bound plus an independent 50-digit verifier. We observe heuristically that the floating-circle cascade continues beyond the certified range, and we state the continuation and the asymptotic form $R^\ast(n)=n^2/8(1+o(1))$ as conjectures. The repository contains the saved certificate artifacts, verifier, and reproducibility commands.
\end{abstract}

\section{Introduction}

The following question was asked by the user Dan on Mathematics Stack Exchange in January 2023 \cite{mse} and appears as question 295 in the UQ benchmark of unsolved Stack Exchange questions \cite{uq}.

\begin{quote}
Non-overlapping circles of radius $1,2,3,\dots,n$ are all externally tangent to a middle circle. How should we arrange the surrounding circles in order to minimize the middle circle's radius $R$?
\end{quote}

Dan conjectured a specific ``pyramid-fashion'' arrangement and observed two phenomena that make the problem subtle: (i) the natural closure equation
\begin{equation}\label{eq:closure}
\sum_{k=1}^{n}\arccos\!\Big(1-\frac{2r_kr_{k+1}}{(R+r_k)(R+r_{k+1})}\Big)=2\pi
\qquad (r_{n+1}=r_1)
\end{equation}
treats consecutive circles as mutually tangent, yet (ii) for $n\ge 8$ the smallest circles can be ``placed last anywhere they fit'', i.e.\ the optimal configuration contains circles tangent to the central circle only. A comment by Rei Henigman \cite{mse} verified the conjecture for $n=10$ by enumerating permutations of $\{2,\dots,10\}$, manually excluding circle $1$ from the chain.

This paper addresses the problem on three levels, keeping the epistemic status of each result explicit.

\paragraph{1. The ordering problem is a Supnick TSP.}
Writing $\theta_R(a,b)$ for the minimum angular separation, at central radius $R$, between the centers of tangent circles of radii $a$ and $b$ (Section~\ref{sec:model}), we show that the matrix $\big(\theta_R(r_i,r_j)\big)_{i,j}$ is symmetric \emph{anti-Monge} for every $R>0$ (Lemma~\ref{lem:supermodular}). Minimizing total angle over cyclic orders is therefore a \emph{maximum} traveling salesman problem on a Supnick matrix, which by a classical theorem of Supnick \cite{supnick} (see also the survey \cite{bdvvw}) is solved by one fixed tour independent of the matrix entries:
\[
\sigma^\ast \;=\; \langle\, 1,\; n-1,\; 3,\; n-3,\; 5,\; n-5,\;\dots,\; n-2,\; 2,\; n \,\rangle .
\]
A short self-consistency argument (the same order is optimal at \emph{every} $R$) transfers optimality from fixed $R$ to the variable-$R$ problem (Theorem~\ref{thm:A}). The order $\sigma^\ast$ coincides with Dan's conjectured pyramid arrangement. In particular $R_{\mathrm{chain}}(\sigma^\ast)$, defined by \eqref{eq:closure}, is an unconditional lower bound for $R^\ast(n)$, and it is the exact optimum whenever the $\sigma^\ast$-necklace is geometrically realizable.

\paragraph{2. Realizability breaks in a finite certified range.}
The closure equation \eqref{eq:closure} ignores non-adjacent pairs. For $n\ge 8$ the configuration it describes for $\sigma^\ast$ is \emph{not realizable}: circle $1$, seated between circles $n$ and $n-1$, no longer keeps them apart, i.e.
\(\theta_R(n,1)+\theta_R(1,n-1) < \theta_R(n,n-1)\).
Circle $1$ then leaves the chain and becomes a \emph{floating} circle, tangent to the central circle only. We model the full problem as a circular system of $\binom n2$ pairwise angular constraints --- equivalently a simple temporal network (STN) --- and give a certified algorithm for it (Section~\ref{sec:algo}). Exhaustive certified optimization for $3\le n\le 14$ (Section~\ref{sec:results}) reveals four regimes, governed by two closed-form quantities of the Supnick tour on $\{2,\dots,n\}$: its geometric validity and the largest Descartes pocket it opens.

\begin{center}
\begin{tabular}{llll}
\toprule
regime & $n$ & floating set & optimal chain \\
\midrule
I (full necklace) & $3\le n\le 7$ & $\varnothing$ & $\sigma^\ast$ on $\{1,\dots,n\}$ \\
II (paid floater) & $n=8,9$ & $\{1\}$ & distorted tour on $\{2,\dots,n\}$ \\
III (free floater) & $n=10,11,12$ & $\{1\}$ & $\sigma^\ast$ on $\{2,\dots,n\}$ \\
IV (second breakdown) & $n=13$ & $\{1\}$ & constrained tour on $\{2,\dots,n\}$ \\
II$'$ (cascade, level 2) & $n=14$ & $\{1,2\}$ & constrained tour on $\{3,\dots,n\}$ \\
\bottomrule
\end{tabular}
\end{center}

\IfFileExists{paper_assets/figures/n14.png}{\begin{figure}[t]\centering\includegraphics[width=.55\textwidth]{paper_assets/figures/n14.png}\caption{The certified optimum for $n=14$; floating circles $1,2$ are highlighted.}\label{fig:n14}\end{figure}}{}

\paragraph{3. Quantitative results.}
We report certified optima, binding structures and floating sets for all $n\le 14$, verified to 50 decimal digits (Table~\ref{tab:main}); the radius-maximizing (worst) arrangement, which is the \emph{minimum}-tour counterpart $\langle 1,3,5,\dots,n,\dots,4,2\rangle$ of Supnick's theorem (Proposition~\ref{prop:worst}); and the conjectured asymptotics $R^\ast(n)\sim n^2/8$ (Section~\ref{sec:asymptotics}).

\paragraph{Epistemic ledger.}
Three levels of certainty are kept explicit throughout. \emph{Proved}: Lemmas~\ref{lem:angular} and~\ref{lem:supermodular}, Theorem~\ref{thm:A}, and the chain-level statement of Proposition~\ref{prop:worst}. \emph{Certified computation} (exhaustive search with explicitly stated numerical guards, independently re-checkable from the repository): Tables~\ref{tab:main} and~\ref{tab:worst}, the realizability statements of Propositions~\ref{prop:break1}--\ref{prop:break3} on their stated finite ranges, and the three seam-failure computations. \emph{Heuristic or conjectural}: Table~\ref{tab:heur}, Conjecture~\ref{conj:cascade}, and Conjecture~\ref{conj:asym}.

\paragraph{Related work.}
<<<<<<< HEAD
Supnick's fixed-tour theorem \cite{supnick} and the Monge/Supnick taxonomy of polynomially solvable TSPs are surveyed in \cite{bdvvw}. Circle packing literature concerns containers (packing circles \emph{inside} regions, see e.g.\ \cite{hifi}); the present problem --- ordering unequal circles \emph{around} a central one --- appears not to have been treated. The tangency mechanics are governed by Descartes' Circle Theorem \cite{mathews-zymaris-spinors}. To our knowledge the link between Supnick matrices and rings of tangent circles is new.
=======
Supnick's fixed-tour theorem \cite{supnick} and the Monge/Supnick taxonomy of polynomially solvable TSPs are surveyed in \cite{bdvvw}. Classical circle-packing theory often fixes tangency or combinatorial data and studies realizations by mutually tangent, non-overlapping circles; Stephenson's monograph is a standard reference for the broader theory \cite{stephenson}. Recent work of Mathews and Zymaris gives Descartes-type algebraic equations for circle-packing configurations, including flower and triangulated-disc settings \cite{mathews-zymaris-spinors,mathews-zymaris-algebraic}. Such algebraic equations could in principle encode configurations related to the rings considered here. The present paper asks a different optimization question: with boundary radii fixed, choose the cyclic order around a central circle and then check all pairwise angular non-overlap constraints. The method therefore exploits the Supnick/anti-Monge ordering structure and STN feasibility checks, rather than solving a prescribed tangency pattern. Packing circles inside containers is a separate optimization literature; see e.g.\ \cite{hifi}. The local tangency computations use Descartes' Circle Theorem \cite{descartes-ref}.
>>>>>>> 2ed20549af73671a647ea86607f0acf2700730eb

\section{The model}\label{sec:model}

Let the central circle have radius $R>0$ and center $O$. A surrounding circle of radius $r$ tangent externally to it has its center at distance $R+r$ from $O$; its position is a single angle $\varphi$. The whole configuration is a vector of angles.

\begin{lemma}[Exact angular reformulation]\label{lem:angular}
Circles of radii $a,b$, both externally tangent to the central circle, with angular positions separated by $\Delta\in[0,\pi]$, are non-overlapping if and only if
\[
\Delta \;\ge\; \theta_R(a,b)\;:=\;2\asin\sqrt{\frac{ab}{(R+a)(R+b)}}.
\]
Moreover $\theta_R(a,b)\in(0,\pi)$ always, $\theta_R$ is symmetric, strictly decreasing in $R$, and strictly increasing in each radius.
\end{lemma}

\begin{proof}
By the law of cosines the squared distance between the centers is $(R+a)^2+(R+b)^2-2(R+a)(R+b)\cos\Delta$, and non-overlap means this is at least $(a+b)^2$. Rearranging gives $\cos\Delta\le 1-\frac{2ab}{(R+a)(R+b)}$, i.e.\ $\Delta\ge\theta_R(a,b)$ with $\theta_R$ as stated, using $\arccos(1-2s^2)=2\asin s$ for $s\in[0,1]$. The argument of the square root is strictly less than $1$ since $(R+a)(R+b)-ab=R(R+a+b)>0$, whence $\theta_R<\pi$; the monotonicity claims are immediate from the formula.
\end{proof}

Thus the problem is: minimize $R$ such that there exist angles $\varphi_1,\dots,\varphi_n$ with circular separations $\ge\theta_R(r_i,r_j)$ for \emph{all} pairs. For a fixed cyclic order this is a circular system of difference constraints --- a simple temporal network --- and feasibility at fixed $R$ is decidable by shortest paths (Section~\ref{sec:algo}). Two derived quantities recur:

\begin{definition}
For a cyclic order $\sigma$ of a multiset of radii, $R_{\mathrm{chain}}(\sigma)$ is the unique solution of $\sum_i \theta_R(\sigma_i,\sigma_{i+1})=2\pi$. (Existence and uniqueness: the left side tends to $|\sigma|\,\pi\ge 2\pi$ as $R\to0^+$, to $0$ as $R\to\infty$, and is strictly decreasing.) For a feasible pair (order, $R$), summing the adjacent gaps shows $\sum_i\theta_R(\sigma_i,\sigma_{i+1})\le2\pi$ and hence $R\ge R_{\mathrm{chain}}(\sigma)$: the chain value is an unconditional per-order lower bound.
\end{definition}

\begin{definition}[Floating circle]
A circle $c$ is \emph{floating} at a global optimum if there exists an optimal configuration in which $c$ is tangent only to the central circle (every pairwise constraint involving $c$ is strictly slack).
\end{definition}

\begin{definition}[Descartes pocket]
For adjacent, mutually tangent circles $a,b$ on the necklace at central radius $R$, the \emph{pocket} is the largest circle externally tangent to the central circle and to both: by Descartes' Circle Theorem its curvature is
$k_{\mathrm{pocket}}=\frac1R+\frac1a+\frac1b+2\sqrt{\frac1{Ra}+\frac1{ab}+\frac1{Rb}}$.
A circle of radius $\rho$ floats freely between $a$ and $b$ iff $\rho\le 1/k_{\mathrm{pocket}}$, equivalently $\theta_R(a,\rho)+\theta_R(\rho,b)\le\theta_R(a,b)$.
\end{definition}

\section{The optimal cyclic order is a fixed Supnick tour}\label{sec:theoremA}

\begin{lemma}[Supermodularity]\label{lem:supermodular}
For every $R>0$ the matrix $\big(\theta_R(r_i,r_j)\big)$, rows and columns indexed by the radii in increasing order, is symmetric strictly anti-Monge: for $a<a'$ and $b<b'$,
\[
\theta_R(a,b)+\theta_R(a',b') \;>\; \theta_R(a,b')+\theta_R(a',b).
\]
\end{lemma}

\begin{proof}
Write $\theta_R(a,b)=2\asin(t_a t_b)$ with $t_x=\sqrt{x/(R+x)}\in(0,1)$, which is strictly increasing in $x$. For $f(x,y)=\asin(xy)$ on $(0,1)^2$,
\[
\frac{\partial^2 f}{\partial x\,\partial y}
=\frac{\partial}{\partial y}\,\frac{y}{\sqrt{1-x^2y^2}}
=\frac{(1-x^2y^2)+x^2y^2}{(1-x^2y^2)^{3/2}}
=(1-x^2y^2)^{-3/2}>0,
\]
so $f$ is strictly supermodular; supermodularity is preserved by strictly increasing reparametrizations of each coordinate, and the claim follows.
\end{proof}

Equivalently, $-\theta_R$ is a Supnick matrix --- symmetric Monge off the diagonal, diagonal entries being irrelevant for tours (cf.\ \cite[Prop.~2.13]{bdvvw}). Supnick's theorem \cite{supnick} covers both extreme tours; in the form stated in \cite{bdvvw} (Theorem~2.5 for the minimum, and the max-TSP discussion following Theorem~2.12, both attributed there to \cite{supnick}), the minimum TSP tour on Supnick matrices is $\langle 1,3,5,\dots,n,\dots,6,4,2\rangle$ and the \emph{maximum} TSP tour is
\[
\sigma^\ast=\langle 1,\; n-1,\; 3,\; n-3,\; 5,\; n-5,\;\dots,\; n-2,\; 2,\; n\rangle,
\]
\emph{independently of the matrix entries}. Maximizing $\sum(-\theta_R)$ is minimizing $\sum\theta_R$, hence:

\begin{theorem}[Optimal order]\label{thm:A}
Let $r_1<\dots<r_n$ be arbitrary positive radii.
\begin{itemize}
\item[(i)] For every $R>0$, among all cyclic orders the tour $\sigma^\ast$ minimizes $\sum_i\theta_R(r_{\sigma(i)},r_{\sigma(i+1)})$.
\item[(ii)] $\sigma^\ast$ minimizes $R_{\mathrm{chain}}$ over all cyclic orders; hence $R^\ast\ge R_{\mathrm{chain}}(\sigma^\ast)$ unconditionally.
\item[(iii)] If the $\sigma^\ast$-necklace at $R_{\mathrm{chain}}(\sigma^\ast)$ is geometrically realizable (no non-adjacent pair violated), then it is a global optimum: $R^\ast=R_{\mathrm{chain}}(\sigma^\ast)$.
\end{itemize}
\end{theorem}

\begin{proof}
(i) is Lemma~\ref{lem:supermodular} plus Supnick's theorem. (ii): fix any order $\sigma$ and let $R=R_{\mathrm{chain}}(\sigma)$. By (i), $\sum\theta_{R}(\sigma^\ast)\le\sum\theta_{R}(\sigma)=2\pi$; since $\sum\theta_R(\sigma^\ast)$ is strictly decreasing in $R$ and equals $2\pi$ at $R_{\mathrm{chain}}(\sigma^\ast)$, we get $R_{\mathrm{chain}}(\sigma^\ast)\le R$. The unconditional bound follows since every feasible $(\sigma,R)$ has $R\ge R_{\mathrm{chain}}(\sigma)$. (iii): realizability makes the lower bound attained.
\end{proof}

\begin{remark}
Theorem~\ref{thm:A} holds for \emph{arbitrary} distinct radii, not only $1,\dots,n$; the arithmetic sequence enters only in the realizability analysis below. For $n=10$ the lower bound of (ii) is $\approx 9.90$, the value observed in \cite{mse} when circle $1$ is (incorrectly) kept in the chain; the true optimum is $9.9799\ldots$.
\end{remark}

\begin{proposition}[Worst arrangement]\label{prop:worst}
For every $R$ the cyclic order maximizing $\sum_i\theta_R$ is the complementary Supnick tour $\sigma_\dagger=\langle1,3,5,\dots,n,\dots,6,4,2\rangle$; consequently $\sigma_\dagger$ maximizes $R_{\mathrm{chain}}$. For all $3\le n\le13$ the $\sigma_\dagger$-necklace is realizable, so among necklace arrangements it maximizes the central radius exactly; certified values in Table~\ref{tab:worst}.
\end{proposition}

\section{Exact algorithm and certification}\label{sec:algo}

\paragraph{Feasibility oracle.} Fix a cyclic order and $R$. Unrolling the circle at one circle ($\varphi_0=0$) gives, for every pair $i<j$ of positions, the two constraints $\varphi_j-\varphi_i\ge\theta_R(r_i,r_j)$ and $\varphi_j-\varphi_i\le2\pi-\theta_R(r_i,r_j)$. This is a simple temporal network; feasibility is the absence of a negative cycle in the associated distance graph and is decided by Floyd--Warshall in $O(n^3)$. The minimum feasible $R$ for the order, $R_{\mathrm{full}}(\sigma)$, is found by bisection (the constraint set relaxes monotonically in $R$), bracketed below by $R_{\mathrm{chain}}(\sigma)$. A witness configuration is recovered from shortest-path potentials and verified independently in Cartesian coordinates.

\paragraph{Lower bounds and search.} For the search over orders we use
\[
\mathrm{LB}(\sigma)=\max\big\{R_{\mathrm{chain}}(\sigma),\,R_{\mathrm{chain}}(\sigma\setminus\{1\}),\,R_{\mathrm{chain}}(\sigma\setminus\{1,2\})\big\},
\]
where $\sigma\setminus S$ denotes the induced cyclic order (valid since induced adjacencies are a subset of the full pairwise constraints). Orders are enumerated canonically (largest radius pinned, reflections halved), $\mathrm{LB}$ is computed vectorized for all $(n-1)!/2$ orders, and $R_{\mathrm{full}}$ is evaluated in ascending $\mathrm{LB}$ until the next bound meets the incumbent: this certifies global optimality up to an absolute guard of $10^{-10}$ in $R$. The guard dominates the floating-point error budget: bisections run to relative tolerance $10^{-13}$; one bound evaluation is an $n$-term sum of $\arcsin$ values, with relative error $O(n\,\varepsilon_{64})$; and the float64 bounds were calibrated against 50-digit recomputation on random samples of orders, with maximal observed deviation $1.8\times10^{-14}$ (over $10^5$ random orders for each $8\le n\le14$; maximal overestimate $1.5\times10^{-14}$), four orders of magnitude below the guard. Independently of the search code, the repository ships a standalone verifier (\texttt{verify.py}) that re-derives each claimed optimum and re-checks every order on the pruning frontier (between $1$ and $11$ orders per $n$) in 50-digit arithmetic. Optimal binding structures were re-verified in 50-digit arithmetic (mpmath); the values in Table~\ref{tab:main} are exact to all displayed digits, with maximal residual slack below $10^{-40}$. The $n=13$ certification enumerates $12!/2\approx2.4\times10^8$ orders ($\approx1.8$\,h on a consumer 8-core laptop); $n=14$ enumerates $13!/2\approx3.1\times10^9$ orders ($\approx22$\,h). Code and artifacts: \url{https://github.com/falker47/ringmin}.

\paragraph{Essential binding structure.} A pairwise constraint is \emph{essential} if perturbing it alone by $+\varepsilon$ makes the STN infeasible at $R^\ast$; a circle is floating iff it carries no essential constraint. This distinguishes true tangencies from placement artifacts of a particular witness.

\section{Certified optima for \texorpdfstring{$n\le 14$}{n <= 14} and the regime cascade}\label{sec:results}

\begin{table}[ht]\centering
\caption{Certified optima, rounded to 12 decimals from the 50-digit verified values (30-digit table in the Appendix). Cycles list radii in circular order; the floating circle is shown in its (non-unique) hosting gap. ``Essential pair'' is the tangent non-adjacent pair bridged by the floater.}
\label{tab:main}
\small
\begin{tabular}{rllll}
\toprule
$n$ & $R^\ast(n)$ & optimal cycle & floating & essential pair \\
\midrule
3 & 0.260869565217 ($=6/23$) & (3,1,2) & --- & --- \\
4 & 0.844453589561 & (4,1,3,2) & --- & --- \\
5 & 1.695494081203 & (5,1,4,3,2) & --- & --- \\
6 & 2.794919518897 & (6,1,5,3,4,2) & --- & --- \\
7 & 4.153189553744 & (7,1,6,3,4,5,2) & --- & --- \\
8 & 5.767794284590 & (8,1,6,4,5,3,7,2) & $\{1\}$ & $(8,6)$ \\
9 & 7.726726552611 & (9,2,8,1,5,6,4,7,3) & $\{1\}$ & $(8,5)$ \\
10 & 9.979907385863 & (10,2,9,4,7,1,6,5,8,3) & $\{1\}$ & $(7,6)$ \\
11 & 12.488720487188 & (11,2,10,4,8,6,7,1,5,9,3) & $\{1\}$ & $(7,5)$ \\
12 & 15.258870430448 & (12,2,11,4,9,6,7,8,5,1,10,3) & $\{1\}$ & $(5,10)$ \\
13 & 18.317563047217 & (13,3,1,12,2,10,6,8,7,9,5,11,4) & $\{1\}$ & $(3,12)$ \\
14 & 21.665395182215 & (14,3,13,2,9,8,7,10,6,11,5,1,12,4) & $\{1,2\}$ & $(13,9),\,(5,12)$ \\
\bottomrule
\end{tabular}
\end{table}

For $n\le7$ the optimal cycle is exactly $\sigma^\ast$ on $\{1,\dots,n\}$ (regime I), as guaranteed by Theorem~\ref{thm:A}(iii) once realizability is checked. The transitions are governed by two computable properties of the Supnick tour on the reduced set $\{2,\dots,n\}$:

\begin{proposition}[First breakdown]\label{prop:break1}
In the $\sigma^\ast$-necklace on $\{1,\dots,n\}$, circle $1$ sits between $n$ and $n-1$. The realizability inequality
$\theta_R(n,1)+\theta_R(1,n-1)\ge\theta_R(n,n-1)$ at $R=R_{\mathrm{chain}}(\sigma^\ast)$
holds for $n\le 7$ and fails for $8\le n\le 13$ (violated pair $(n,n-1)$; at $n=13$ also $(13,11)$ across circle $2$). Hence regime I ends at $n=7$. We conjecture failure for all $n\ge8$; cf.\ open problem (1).
\end{proposition}

\begin{proposition}[Free vs.\ paid floater]\label{prop:break2}
Let $\rho_{\max}(n)$ be the largest Descartes pocket of the $\sigma^\ast$-necklace on $\{2,\dots,n\}$ at its chain radius. Then $\rho_{\max}=0.858,\,0.988,\,1.139,\,1.262,\,1.409$ for $n=8,\dots,12$. For $n=10,11,12$ one has $\rho_{\max}>1$: circle $1$ floats freely and
$R^\ast(n)=R_{\mathrm{chain}}\big(\sigma^\ast\text{ on }\{2,\dots,n\}\big)$ exactly (regime III).
For $n=8,9$, $\rho_{\max}<1$: the optimum distorts the reduced tour to open a unit pocket, at certified extra cost $\delta(8)=0.04080$, $\delta(9)=0.00532$ over the (then unattainable) reduced chain bound (regime II).
\end{proposition}

\begin{proposition}[Second breakdown]\label{prop:break3}
At $n=13$ the Supnick tour on $\{2,\dots,13\}$ is itself unrealizable: circle $2$, seated between $13$ and $12$, violates
$\theta_R(13,2)+\theta_R(2,12)=0.80448<0.83488=\theta_R(13,12)$ at $R=18.3176$.
The certified optimum demotes circle $2$ to the pair $(12,10)$, where the separation constraint holds with slack $6\times10^{-3}$ rad; circle $1$ floats in the $(3,12)$ pocket (radius $1.086$). Thus the mechanism that expelled circle $1$ at $n=8$ reaches circle $2$ at $n=13$ (regime IV).
\end{proposition}

\begin{remark}[Structural summary]
For all $n\le14$ the optimum has the form: a single tangent necklace carrying a Supnick-like order \emph{constrained to geometric realizability}, plus floating circles in Descartes pockets. The cascade continues self-similarly: circle $2$ first floats at $n=14$ (now \emph{certified}, Table~\ref{tab:main}) via a \emph{paid} distortion that opens a pocket of radius $2.013$ between circles $13$ and $9$ --- the regime-II mechanism one level up, at extra cost $0.0565$ over the (then unattainable) chain bound of the Supnick tour on $\{3,\dots,14\}$. At $n=16$ the best known value coincides to $10^{-10}$ with the chain value of the Supnick tour on $\{3,\dots,16\}$, the level-2 analogue of regime III. Heuristically, circle $3$ joins by $n=18$ (Table~\ref{tab:heur}). Moreover the Supnick tour on $\{3,\dots,n\}$ itself becomes unrealizable at $n=17$ --- circle $3$ failing at the $(17,16)$ seam with slack $-0.0104$ --- so the \emph{seam-failure} onsets for circles $1,2,3$ are $n=8,13,17$: each \emph{computed} cascade level reproduces the full regime pattern of the previous one. The float onsets $n=8,14,18$ for circles $1,2,3$ are consistent with Dan's $n/4$ Descartes-based upper bound on the number of ignorable circles \cite{mse}. Whether the cascade continues indefinitely is the content of the following conjecture.
\end{remark}

\begin{conjecture}[Cascade]\label{conj:cascade}
For every $k\ge1$ there is a threshold $s_k$ such that the Supnick tour on $\{k,\dots,n\}$ is geometrically unrealizable for all $n\ge s_k$ (observed and certified: $s_1=8$, $s_2=13$, $s_3=17$), and circle $k$ is floating in every optimal configuration for all sufficiently large $n$, with each level exhibiting the paid-then-free regime pattern of Propositions~\ref{prop:break1}--\ref{prop:break3}.
\end{conjecture}

\begin{table}[ht]\centering
\caption{Best known arrangements for $15\le n\le18$ (local search, 200--600 restarts per $n$; \emph{non-exhaustive}). The last column counts restarts reaching the best value.}
\label{tab:heur}
\small
\begin{tabular}{rlllr}
\toprule
$n$ & best $R$ & floating & cycle & hits \\
\midrule
15 & 25.290151636714 & $\{1,2\}$ & (15,3,14,5,12,2,9,1,8,10,7,11,6,13,4) & 59 \\
16 & 29.194988498141 & $\{1,2\}$ & (16,1,4,14,6,12,8,10,9,2,11,7,13,5,15,3) & 18 \\
17 & 33.394425022062 & $\{1,2\}$ & (17,4,16,3,15,6,1,13,8,11,10,9,2,12,7,14,5) & 19 \\
18 & 37.870755667393 & $\{1,2,3\}$ & (18,4,17,3,14,8,12,10,11,9,2,13,7,15,6,16,1,5) & 9 \\
\bottomrule
\end{tabular}
\end{table}

\section{Asymptotics}\label{sec:asymptotics}

For $R\gg n$ one has $\theta_R(a,b)=2\sqrt{ab}/R\,\big(1+O(n/R)\big)$, so the closure equation gives
$R\approx\frac1\pi\sum_i\sqrt{r_ir_{i+1}}$. Since $\sqrt{ab}$ is supermodular, Theorem~\ref{thm:A}(i) applies verbatim: the minimizing order of $\sum\sqrt{r_ir_{i+1}}$ is again $\sigma^\ast$, whose value on $1,\dots,n$ is
\[
\min_\sigma\sum_i\sqrt{r_{\sigma(i)}r_{\sigma(i+1)}}
=\Big(1+o(1)\Big)\int_0^n\sqrt{x(n-x)}\,dx=\frac{\pi n^2}{8}\,(1+o(1)),
\]
because adjacent pairs in $\sigma^\ast$ are $\{k,\,n-k+O(1)\}$ with each odd $k$ contributing twice. Self-consistency ($R\sim n^2/8\gg n$) further requires the floating circles to be asymptotically negligible, in the sense $\sum_{k\in F(n)}\sqrt{kn}=o(n^2)$ --- guaranteed if $|F(n)|=o(n)$, consistent with $|F(18)|=3$, but open (problem (2) below). Under this assumption we are led to:

\IfFileExists{paper_assets/figures/radii_vs_n.png}{\begin{figure}[t]\centering\includegraphics[width=.78\textwidth]{paper_assets/figures/radii_vs_n.png}\caption{$R^\ast(n)$, the worst-arrangement curve, and the $n^2/8$ asymptote.}\label{fig:radii-vs-n}\end{figure}}{}

\begin{conjecture}\label{conj:asym}
$\displaystyle R^\ast(n)=\frac{n^2}{8}\big(1+o(1)\big)$. Empirically the deficit $n^2/8-R(n)$ equals $2.52,\,2.64,\,2.74,\,2.81,\,2.83$ for the certified $n=10,\dots,14$, and $2.83,\,2.81,\,2.73,\,2.63$ for the heuristic $n=15,\dots,18$ (the latter are \emph{lower} bounds on the true deficit, heuristic radii being upper bounds for $R^\ast$): bounded so far, and in any case growing strictly slower than $n$. We tentatively conjecture the stronger form $n^2/8-R^\ast(n)=O(\sqrt n\,)$.
\end{conjecture}

Two-sided bounds making the leading order rigorous (via $\asin x\ge x$ and the explicit $\sigma^\ast$ configuration as an upper bound) appear within reach but are not pursued here.

\section{Open problems}

(1) Prove Propositions~\ref{prop:break1}--\ref{prop:break3} for all $n$ analytically (the inequalities are explicit; monotonicity in $n$ is the missing step). (2) Characterize the floating set $F(n)$ asymptotically (float onsets: circle $1$ at $n=8$ and circle $2$ at $n=14$, both certified, circle $3$ at $n=18$ heuristically; certified seam-failure onsets $n=8,13,17$); is $|F(n)|\sim cn^{1/2}$, or does Dan's $n/4$ bound give the right order? Both are subsumed by proving Conjecture~\ref{conj:cascade}. (3) The same machinery applies to radii $r_k=k^\alpha$ or arbitrary sequences: Theorem~\ref{thm:A} is general, the cascade analysis is not. (4) Three dimensions: spheres of radii $1,\dots,n$ tangent to a central sphere --- the ordering structure disappears; what replaces Supnick? 

\section*{Acknowledgments}
The problem and the pyramid conjecture are due to Dan \cite{mse}; Rei Henigman's $n=10$ computation provided an independent check. The project used AI assistance for software, experimentation, drafting, and repository organization under the author's direction and final review. The paper is written to be independently checkable: the proof of Theorem~\ref{thm:A} is short and self-contained, the only imported result is Supnick's theorem, and every certified value is reproducible from the public repository with the independent verifier described in Section~\ref{sec:algo}.

\begin{thebibliography}{9}
\bibitem{mse} Dan, \emph{Circles of radius $1,2,3,\dots,n$ all touch a middle circle. How to make the middle circle as small as possible?}, Mathematics Stack Exchange, question 4619480, January 2023. \url{https://math.stackexchange.com/q/4619480}
\bibitem{uq} The UQ Project, \emph{Unsolved Questions} benchmark, question 295. \url{https://uq.stanford.edu/question/295} (accessed June 2026).
\bibitem{supnick} F. Supnick, \emph{Extreme Hamiltonian lines}, Annals of Mathematics (2) \textbf{66} (1957), 179--201.
\bibitem{bdvvw} R.E. Burkard, V.G. De\u{\i}neko, R. van Dal, J.A.A. van der Veen, G.J. Woeginger, \emph{Well-solvable special cases of the TSP: a survey}, SIAM Review \textbf{40}(3) (1998), 496--546.
<<<<<<< HEAD
\bibitem{mathews-zymaris-spinors} D.V. Mathews, O. Zymaris, \emph{Spinors and the Descartes circle theorem}, Journal of Geometry and Physics \textbf{212} (2025), 105458, doi:10.1016/j.geomphys.2025.105458.
=======
\bibitem{stephenson} K. Stephenson, \emph{Introduction to Circle Packing: The Theory of Discrete Analytic Functions}, Cambridge University Press, 2005.
\bibitem{mathews-zymaris-spinors} D.V. Mathews, O. Zymaris, \emph{Spinors and Descartes' Theorem}, arXiv:2310.11701, 2023.
\bibitem{mathews-zymaris-algebraic} D.V. Mathews, O. Zymaris, \emph{An algebraic approach to circle packing}, arXiv:2504.14593, 2025.
\bibitem{descartes-ref} J.C. Lagarias, C.L. Mallows, A.R. Wilks, \emph{Beyond the Descartes circle theorem}, Amer. Math. Monthly \textbf{109} (2002), 338--361.
>>>>>>> 2ed20549af73671a647ea86607f0acf2700730eb
\bibitem{hifi} M. Hifi, R. M'Hallah, \emph{A literature review on circle and sphere packing problems: models and methodologies}, Advances in Operations Research (2009).
\end{thebibliography}

\appendix
\section{High-precision values and figures}
\begin{table}[ht]\centering
\caption{The radius-maximizing necklace arrangement $\sigma_\dagger$ (Proposition~\ref{prop:worst}): certified chain values, all geometrically realizable.}
\label{tab:worst}
\small
\begin{tabular}{lrrrrrrrrrrr}
\toprule
$n$ & 3 & 4 & 5 & 6 & 7 & 8 & 9 & 10 & 11 & 12 & 13 \\
$R_\dagger$ & 0.2609 & 0.8673 & 1.7786 & 3.0001 & 4.5375 & 6.3934 & 8.5690 & 11.0647 & 13.8806 & 17.0164 & 20.4720 \\
\bottomrule
\end{tabular}
\end{table}

\IfFileExists{paper_assets/appendix_tables.tex}{\begin{table}
\centering
\begin{tabular}{rl}
\toprule
$n$ & $R^*$ \\
\midrule
3 & 0.260869565217391304347826086957 \\
4 & 0.844453589560855604347528524674 \\
5 & 1.69549408120271081351328017371 \\
6 & 2.79491951889692485617024406797 \\
7 & 4.15318955374381246513863858202 \\
8 & 5.7677942845896143026361805725 \\
9 & 7.72672655261128921886246177604 \\
10 & 9.97990738586347760966552641468 \\
11 & 12.4887204871876588517468786264 \\
12 & 15.2588704304484933617250043503 \\
13 & 18.3175630472173206282821941532 \\
\bottomrule
\end{tabular}
\caption{Certified optima, 30 significant digits.}
\end{table}

\begin{table}
\centering
\begin{tabular}{rlll}
\toprule
$n$ & Supnick min tour & $R_{\mathrm{chain}}$ & $R_{\mathrm{full}}$ \\
\midrule
3 & [1, 3, 2] & 0.260869565217 & 0.260869565217 \\
4 & [1, 3, 4, 2] & 0.867281640255 & 0.867281640255 \\
5 & [1, 3, 5, 4, 2] & 1.778550949884 & 1.778550949884 \\
6 & [1, 3, 5, 6, 4, 2] & 3.000138442323 & 3.000138442323 \\
7 & [1, 3, 5, 7, 6, 4, 2] & 4.537467833562 & 4.537467833562 \\
8 & [1, 3, 5, 7, 8, 6, 4, 2] & 6.393358320114 & 6.393358320114 \\
9 & [1, 3, 5, 7, 9, 8, 6, 4, 2] & 8.568988935955 & 8.568988935955 \\
10 & [1, 3, 5, 7, 9, 10, 8, 6, 4, 2] & 11.064726494858 & 11.064726494858 \\
11 & [1, 3, 5, 7, 9, 11, 10, 8, 6, 4, 2] & 13.880577932626 & 13.880577932626 \\
12 & [1, 3, 5, 7, 9, 11, 12, 10, 8, 6, 4, 2] & 17.016408238128 & 17.016408238128 \\
13 & [1, 3, 5, 7, 9, 11, 13, 12, 10, 8, 6, 4, 2] & 20.472039120681 & 20.472039120681 \\
14 & [1, 3, 5, 7, 9, 11, 13, 14, 12, 10, 8, 6, 4, 2] & 24.247291190092 & 24.247291190092 \\
15 & [1, 3, 5, 7, 9, 11, 13, 15, 14, 12, 10, 8, 6, 4, 2] & 28.342000288518 & 28.342000288518 \\
16 & [1, 3, 5, 7, 9, 11, 13, 15, 16, 14, 12, 10, 8, 6, 4, 2] & 32.756022334220 & 32.756022334220 \\
17 & [1, 3, 5, 7, 9, 11, 13, 15, 17, 16, 14, 12, 10, 8, 6, 4, 2] & 37.489233291550 & 37.489233291550 \\
18 & [1, 3, 5, 7, 9, 11, 13, 15, 17, 18, 16, 14, 12, 10, 8, 6, 4, 2] & 42.541527273519 & 42.541527273519 \\
\bottomrule
\end{tabular}
\caption{Worst-arrangement comparison; Supnick min tour values.}
\end{table}
}{\emph{[30-digit tables: compile with \texttt{paper\_assets/} present, or see the repository.]}}

Per-$n$ configuration figures for all $3\le n\le 18$ are available in the repository.

\end{document}
